\section{Introduction}

% add paragraphs for introduction

Presentation slides: \url{https://leeds365-my.sharepoint.com/:p:/g/personal/fy22az_leeds_ac_uk/IQAzCfArJhdOSoHPXOb2jwl4Acw6uB58h-HJmVDQevXYUGE?e=V1H29b}

\subsection{Dataset}

Global - Food Prices

\url{https://data.humdata.org/dataset/global-wfp-food-prices}

% quote this
This dataset contains Countries, Commodities and Markets data, sourced from the World Food Programme Price Database. The World Food Programme Price Database covers foods such as maize, rice, beans, fish, and sugar for 98 countries and some 3000 markets. It is updated weekly but contains to a large extent monthly data. The data goes back as far as 1992 for a few countries, although many countries started reporting from 2003 or thereafter.

License
Creative Commons Attribution for Intergovernmental Organisations (CC BY-IGO)
\url{https://creativecommons.org/licenses/by/3.0/igo/legalcode}

% Discuss allowed for this project.

\section{Design}

\subsection{Features}

% CRUD features
% User System
% Data Explorer – CRUD Operations on prices

% Advanced features - Analytics
Price trends – in a specific country
Commodity volatility – using Coefficient of Variation (CV)
Regional comparison – which regions face high food costs
Market summary – number of prices, commodities, latest prices
Crisis scores – food-crisis indicator (0-100)

%  Advanced features - MCP Server
get_global_crisis_overview – top 20 most crisis countries
get_crisis_summary – same as before, for a specific country
get_price_trends – same
compare_regional_prices – same
get_volatile_commodities – same

% Advanced Features - Interactive choropleth map

% Frontend
% screenshot frontend.png

World Map, Explorer, Data Entry
API Documentation - Swagger UI, ReDOC

\subsection{System Architecture}

two figures in figures/, architecture.svg and  ERD.html

% write in paragrams, give more reasons and critical analysis of the design choices
Database
PostgreSQL – good performance, best practice
SQLAlchemy – Object Relational Mapping (ORM), not need write SQL
Alembic – migration
Backend
Python – Pythonic, concise, easy to read/write
FastAPI – modern framework, less code, generated API doc
Pydantic – verification, typing system
Frontend
React – component, reactive, Single Page Application (SPA)
Typescript – type system
Tailwind CSS – common practice, easy for GenAI
Leaflet.js – map
Chat.js – charts
Deployment
Docker (Compose) – easy to migrate
Testing
Pytest, pytest-cov – pytest and coverage
FastAPI TestClient, httpx – http endpoints

\subsection{Database Design}

% describe db design, make relation to dataset

\subsection{API Design}

% list instructions

API Documentation: \url{https://github.com/COMP3011-Coursework/COMP3011-Coursework-1/blob/main/docs/api-reference.pdf}

\section{Implementation}

GenAI chat for
Explore project idea, dataset
Decide tech stack
Discuss initial design
Manually write an initial design document
Iterate this document with GenAI
Use Claude Code to coding with this given design document
Start project and check design, functions, issues
Fix and iterate project with Claude Code

\section{Testing}

Automated testing frameworks – pytest, FastAPI TestClient
100% coverage
Manually frontend test – design, features
MCP tests – using Claude Code connected to MCP server
Ask friends to test

\section{Deployment}

Support
Self-hosted with Docker Compose
Render.com \url{https://render.com/}
Live Demo
Frontend \url{https://comp3011-coursework-1-frontend.onrender.com/}
Backend — API docs (Swagger) \url{https://comp3011-coursework-1-backend.onrender.com/docs}
Backend — API docs (ReDoc) \url{https://comp3011-coursework-1-backend.onrender.com/redoc}

\section{Performance}

Local environment
~3.14M prices
~50ms page loading time
< 1s loading world score loading time and ~100ms data loading time
Live demo
~540k (year <=2015) prices
<200ms page loading time
~5s and ~2s time corresponding
Cached GeoJSON — ~14MB, cached at client
Pagination used, page size parameter

\section{Reflections}

\subsection{Version Control Practices}

Git, 60+ commits
Self-hosted Gitea with private repository
Public GitHub repository released before presentation
Reason for not to use branches for this project
Simple
No cooperation or external customer
Agile
Conventional Commits \url{https://www.conventionalcommits.org/en/v1.0.0/} used
A practice to write commit messages
feat, fix, docs, perf, test, build, ...

\subsection{GenAI Usage}

Project Idea
Explore Datasets and Understanding Format
Structure Design and Tech Stack
Agentic Development and Vibe Coding
Debugging and Fixing
Generate Diagrams
Review the report and presentation

\section{Limitations and Future Work}

% discuss limitations and future work

\section{Conclusions}

% write conclusions

\subsection{Deliverables}

Deployed website
https://comp3011-coursework-1-frontend.onrender.com
GitHub code repository
https://github.com/COMP3011-Coursework/COMP3011-Coursework-1
API documentation
https://comp3011-coursework-1-backend.onrender.com/docs
https://comp3011-coursework-1-backend.onrender.com/redoc
PDF link available in technical report
Technical report - submitted via Minerva
Presentation slides - submitted via Minerva

GitHub repository: \url{https://github.com/COMP3011-Coursework/COMP3011-Coursework-1/tree/main}
